\documentclass[dvipsnames]{beamer}
\usepackage{comment}
\usepackage{hyperref}

\title{MATH 1190 \& 1210 Meeting for Week 13}
\author{Josh}
\date{6 April 2026}

\newcommand{\transition}[1]{\frame{{}\begin{center}\fbox{#1}\end{center}}}

\begin{document}

\maketitle


\section{Logistics: Pacing}

\transition{Logistics}

\frame{{Logistics: Pacing}

{\bf Where are folks in the lessons right now?}

}

\frame[t]{{Logistics: Upcoming Schedule}

    \begin{itemize}
     \item week 13 of class (6 April - 10 April)
        \begin{itemize}
        \item lesson 20 (Discrete RVs)
        \item Application Day 3: Gradient Descent
        \end{itemize}
         \item week 14 of class (13 April - 17 April)
        \begin{itemize}
        \item lesson 21 (Continuous RVs)
        \item lesson 22 (Expectation and Variance) -- {\footnotesize\color{red} last topic on assessment 4}
        \end{itemize}
    \end{itemize}

}

\frame[t]{{Logistics: Upcoming Deadlines}

{\bf Assessment-related Logistics}
\begin{itemize}
\item KC2 Reflections
\item The KC3 Grading (done by end of day Thursday 9 April)
\item KC2 Regrade requests some still need finished F1 (not on KC3)
\item KC3 Regrade requests open on Saturday 11 April due by Thursday 16 April
\item KC4 Week after next 23 April
    \begin{itemize}
    \item reassessment targets: {\color{RedOrange}A2},{\color{RedOrange}A3}, {\color{Fuchsia}I1}, {\color{Fuchsia}I2}, 
    \item new targets:{\color{Fuchsia}I3},{\color{Fuchsia}I4},{\color{ProcessBlue}F3}, {\color{Orange}P1}, {\color{Orange}P2} (5 targets!)
    \item Thursday, 23 April, 7:00 PM; 80 minutes
    \item Room assignments same as KC1 this time. No switching dates!
    \end{itemize}
\end{itemize}

}




\section{Content}

\transition{Content}


%------------------------------------------------

%------------------------------------------------

\begin{frame}{Learning Objectives: P1: Introduction to Probability }
After this lesson, students will be able to:
\begin{itemize}
    \item Describe sample spaces and events
    \item Identify discrete random variables
    \item Compute probabilities from finite sample spaces
    \item Understand and use probability mass functions (PMFs)
    \item Interpret and sketch histograms of PMFs
\end{itemize}
\end{frame}

%------------------------------------------------

\begin{frame}{Basic Terminology}
\begin{itemize}
    \item \textbf{Experiment}: A repeatable process with uncertain outcomes
    \item \textbf{Outcome}: A single result of an experiment
    \item \textbf{Sample Space ($S$)}: Set of all possible outcomes
    \item \textbf{Event}: A subset of the sample space
\end{itemize}
\end{frame}


%------------------------------------------------

\begin{frame}{Discrete Random Variables}
\begin{itemize}
    \item Warmup: Exercise A: Drawing Marbles
    \item \textbf{Definition of a random variable} assigns a number to each outcome
    \item \textbf{Difference between discrete and continuous} takes on a finite number of values
    \item Definition of uniform probability calculation
    \item Definition of histogram
\end{itemize}
\end{frame}


\begin{frame}{Discrete Random Variables}

\textbf{Example 1:Random variable on sample space of Exercise A}
\begin{itemize}
    \item Create histogram
    \item Build intuition for random variables
\end{itemize}
\end{frame}

%------------------------------------------------

\begin{frame}{Properties of Probability Mass Function}

\begin{itemize}
    \item $0 \le p(x) \le 1$
    \item $\sum p(x) = 1$
\end{itemize}

\textbf{Exercise B}
\begin{itemize}
    \item Fill in table of values
    \item Create histogram
\end{itemize}
\end{frame}

%------------------------------------------------

%------------------------------------------------

%------------------------------------------------

\begin{frame}{Histograms of PMFs}
\begin{itemize}
    \item Exercise C: Histogram visually represents a PMF
    \item $x$-axis: possible values
    \item $y$-axis: probabilities
\end{itemize}



\begin{block}{Interpretation}
Height of each bar = $P(X = x)$
\end{block}
\end{frame}

%------------------------------------------------

\begin{frame}{Exercise D: Classic two dice example}
\begin{itemize}
    \item Let $X =$ sum of two dice
    \item Possible values: $2,3,\dots,12$
\end{itemize}

\textbf{Key idea:}
\begin{itemize}
    \item Not all sums are equally likely!
    \item Example:
    \[
    P(X=7) = \frac{6}{36}, \quad P(X=2)=\frac{1}{36}
    \]
\end{itemize}
\end{frame}

%------------------------------------------------

\begin{frame}{Applications: Age Distribution}
\begin{itemize}
    \item Discrete random variable can model real data
    \item Example: Age groups of a population
\end{itemize}


\begin{block}{Key Idea}
Probabilities represent proportions of a population
\end{block}

\begin{itemize}
    \item Histograms help visualize distributions
    \item Area corresponds to total probability
\end{itemize}
\end{frame}

%------------------------------------------------

\begin{frame}{Summary}
\begin{itemize}
    \item Sample space = all possible outcomes
    \item Events = subsets of outcomes
    \item Discrete random variables take finite values
    \item PMF describes probabilities
    \item Histograms visualize distributions
\end{itemize}
\end{frame}

%------------------------------------------------

%%%%%%%%%%%%%%%%%%%%%%%%%%%%%%%%%%%%%%%%%%%%%%%%%%%%%%%%%%%%%%%%%%%%%%%%%%%%%%%%%%%%%%%%%%%%%%%%%%%%%%
%%%%%%%%%%%%%%%%%%%%%%%%%%%%%%%%%%%%%%%%%%%%%%%%%%%%%%%%%%%%%%%%%%%%%%%%%%%%%%%%%%%%%%%%%%%%%%%%%%%%%%




%------------------------------------------------

%------------------------------------------------

\begin{frame}{Application Day 3: Gradient Descent}
\begin{itemize}
    \item Machine learning = optimization problem
    \item We define a one variable loss function measuring error
    \item Goal: Use basic calculus ideas to minimize loss when we cannot solve using algebra or exact derivatives
\end{itemize}
\end{frame}

%------------------------------------------------
\section{Pre-Class Activity}

\begin{frame}{Pre-Class: Key Idea}
\begin{itemize}
    \item Use derivatives to guide optimization
    \item Move in direction of \textbf{steepest decrease}
\end{itemize}



\[
x_{n+1} = x_n - \alpha f'(x_n)
\]



\begin{itemize}
    \item Use a hyperparameter $\alpha$ = learning rate (step size)
    \item Repeat until close to minimum
\end{itemize}

\end{frame}

%------------------------------------------------

\begin{frame}{Pre-Class: Intuition}
\begin{itemize}
    \item If $f'(x) > 0$ $\rightarrow$ move left
    \item If $f'(x) < 0$ $\rightarrow$ move right
    \item If $f'(x) \approx 0$ $\rightarrow$ near minimum
\end{itemize}



\begin{block}{Tradeoff}
\begin{itemize}
    \item Large $\alpha$ $\rightarrow$ may overshoot
    \item Small $\alpha$ $\rightarrow$ slow progress
\end{itemize}
\end{block}

\end{frame}

%------------------------------------------------

\begin{frame}{Pre-Class: Example Process}
\begin{itemize}
    \item Start with $x_1$
    \item Compute $f'(x_1)$
    \item Update:
    \[
    x_2 = x_1 - \alpha f'(x_1)
    \]
 %   \item Repeat to get $x_3, x_4, \dots$
\end{itemize}



\begin{block}{Goal}
Sequence approaches a minimum of $f(x)$
\end{block}

\end{frame}

%------------------------------------------------

\begin{frame}{Pre-Class: Exploration}
Using the app:
\begin{itemize}
    \item Try different functions
    \item Change starting value $x_0$
    \item Adjust learning rate $\alpha$
\end{itemize}



\textbf{Questions to consider:}
\begin{itemize}
    \item Does it always find the minimum?
    \item What happens if $\alpha$ is large?
    \item How does the starting point matter?
\end{itemize}

\end{frame}

%------------------------------------------------
\section{During-Class Activity}

\begin{frame}{During Class: Real-World Application}
\begin{itemize}
    \item Model spread of disease
    \item Early growth is approximately exponential:
\end{itemize}

\[
I(t) = e^{\lambda t}
\]



\begin{itemize}
    \item Take logarithm:
\end{itemize}

\[
\ln I(t) = \lambda t
\]



\begin{block}{Goal}
Estimate $\lambda$ using gradient descent
\end{block}

\end{frame}

%------------------------------------------------

\begin{frame}{During Class: Error Function}
We measure how well our model fits data:

\[
E(\lambda)
\]


\begin{itemize}
    \item $y_i = \ln I(t_i)$
    \item Minimize error using gradient descent
\end{itemize}

\end{frame}

%------------------------------------------------

\begin{frame}{During Class: Gradient Descent Steps}
\begin{enumerate}
    \item Choose initial guess $\lambda_0$
    \item Compute predictions $\hat{y}_i = \lambda t_i$
    \item Compute derivative:
    \[
    \frac{dE}{d\lambda} 
    \]
    \item Update:
    \[
    \lambda_{\text{new}} = \lambda_{\text{old}} - \alpha \frac{dE}{d\lambda}
    \]
\end{enumerate}

\end{frame}

%------------------------------------------------

\begin{frame}{During Class: Interpretation}
Once we estimate $\lambda$, we compute:

$R_0$  = ``how many people, on average, one infected person will infect." 




\begin{itemize}
    \item $R_0 > 1$ $\rightarrow$ disease spreads
    \item $R_0 = 1$ $\rightarrow$ stable
    \item $R_0 < 1$ $\rightarrow$ dies out
\end{itemize}

\end{frame}

%------------------------------------------------

\begin{frame}{During Class: Exploration}
Using the app:
\begin{itemize}
    \item Run gradient descent on real data
    \item Record $\lambda$, error, and $R_0$
\end{itemize}



\textbf{Investigate:}
\begin{itemize}
    \item Effect of learning rate $\alpha$
    \item Effect of number of steps
    \item Best parameter choices
    \item Calculate $R_0$ 
    \item Suggest course of action to epidimiologist
\end{itemize}

\end{frame}

%------------------------------------------------



%------------------------------------------------


%%%%%%%%%%%%%%%%%%%%%%%%%%%%%%%%%%%%%%%%%%%%%%%%%%%%%%%%%%%%%%%%%%%%%%%%%%%%%%%%%%%%%%%%%%%%%%%%%%%%%%
%%%%%%%%%%%%%%%%%%%%%%%%%%%%%%%%%%%%%%%%%%%%%%%%%%%%%%%%%%%%%%%%%%%%%%%%%%%%%%%%%%%%%%%%%%%%%%%%%%%%%%

\section{Pedagogy}

\transition{Pedagogy}


\transition{Active Learning}
\transition{And Now - An Active Learning Discussion with Jim B.!}
%%%%%%%%%%%%%%%%%%%%%%%%%%%%%%%%%%%%%%%%%%%%%%%%%%%%%%%%%%%%%%%%%%%%%%%%%%%%%%%%%%%%%%%%%%%%%%%%%%%%%%

%------------------------------------------------



%------------------------------------------------
\section{Motivation}

\begin{frame}{Today's Paper}
``Interventions to reduce burnout in students: A systematic
review and meta‑analysis" Madigan et al. (2024)

\begin{itemize}
    \item Student burnout is increasing in higher education
    \item Linked to:
    \begin{itemize}
        \item Lower academic performance
        \item Reduced engagement
        \item Mental health challenges
    \end{itemize}
\end{itemize}



\begin{block}{Key Question}
What interventions actually reduce burnout in students?
\end{block}

\end{frame}

%------------------------------------------------
\section{What is Burnout?}

\begin{frame}{Defining Student Burnout}
Burnout consists of three dimensions:
\begin{itemize}
    \item \textbf{Emotional exhaustion}
    \item \textbf{Cynicism / disengagement}
    \item \textbf{Reduced academic efficacy}
\end{itemize}



\begin{block}{Important Distinction}
Burnout is not just stress — it is chronic, cumulative, and affects identity as a student
\end{block}

\end{frame}

%------------------------------------------------
\section{Study Overview}

\begin{frame}{Study Design}
\begin{itemize}
    \item Type: Systematic review + meta-analysis
    \item Synthesizes many intervention studies
    \item Focus: What works to reduce student burnout?
\end{itemize}



\begin{block}{Goal}
Estimate the effectiveness of different intervention types
\end{block}

\end{frame}

%------------------------------------------------

\begin{frame}{Data and Methods}
\begin{itemize}
    \item Collected studies from multiple academic databases

    \begin{itemize}
    \item The authors did not recruit students
    \item They did not run interventions themselves
    \item Experimental or quasi-experimental designs
    \item Measured burnout outcomes
    \end{itemize}
    \item Compared:
    \begin{itemize}
        \item Intervention vs. control groups
        \item paper asks “Across all studies, how well do interventions work?”
    \end{itemize}
\end{itemize}


\begin{block}{Analysis}
Meta-analysis to compute overall effect sizes
\end{block}

\end{frame}


%------------------------------------------------
\section{Measuring Burnout}

\begin{frame}{How is Burnout Measured?}
\begin{itemize}
    \item Burnout is not directly observable
    \item Measured using validated survey instruments
    \item Students respond to statements using a Likert scale
    \begin{itemize}
        \item Example: 1 = Never, 5 = Always
    \end{itemize}
\end{itemize}
\end{frame}

%------------------------------------------------

\begin{frame}{Common Burnout Scales}
\textbf{Maslach Burnout Inventory -- Student Survey (MBI-SS)}
\begin{itemize}
    \item Emotional exhaustion
    \item Cynicism / disengagement
    \item Academic efficacy (reverse-coded)
\end{itemize}



\textbf{Oldenburg Burnout Inventory (Student Version)}
\begin{itemize}
    \item Exhaustion
    \item Disengagement
\end{itemize}





\end{frame}

%------------------------------------------------

\begin{frame}{Example Survey Items}
Students rate statements such as:

\begin{itemize}
    \item ``I feel emotionally drained by my studies''
    \item ``I have become less interested in my coursework''
    \item ``I feel capable of succeeding in my classes''
\end{itemize}



\begin{block}{Interpretation}
Higher exhaustion + disengagement + lower efficacy = higher burnout
\end{block}

\end{frame}

%------------------------------------------------

\begin{frame}{From Surveys to Data}
\begin{itemize}
    \item Each student receives a burnout score
    \item Studies compare:
    \begin{itemize}
        \item Intervention group vs. control group
        \item Before vs. after intervention
    \end{itemize}
\end{itemize}



\begin{itemize}
    \item Results converted into effect sizes
\end{itemize}




\end{frame}
%------------------------------------------------
\section{Key Results}

\begin{frame}{Main Findings}
\begin{itemize}
    \item Interventions \textbf{do reduce burnout overall}
    \item Effects are:
    \begin{itemize}
        \item Small to moderate
        \item Statistically significant
    \end{itemize}
\end{itemize}





\end{frame}
%------------------------------------------------
\section{Intervention Strategies}

\begin{frame}{Examples of Burnout Interventions}
Studies in the paper implemented interventions such as:



\textbf{Cognitive / Behavioral}
\begin{itemize}
    \item Mindfulness and meditation exercises
    \item Cognitive Behavioral Therapy (CBT) techniques
    \item Stress-reframing strategies
\end{itemize}



\textbf{Academic Skills}
\begin{itemize}
    \item Time management training
    \item Study skills workshops
    \item Goal-setting and planning exercises
\end{itemize}
\end{frame}
\begin{frame}


\textbf{Wellness and Physical Activity}
\begin{itemize}
    \item Exercise programs
    \item Relaxation techniques (breathing, stretching)
\end{itemize}



\textbf{Social Support}
\begin{itemize}
    \item Peer discussion groups
    \item Counseling or mentoring programs
\end{itemize}

\end{frame}






%------------------------------------------------



%------------------------------------------------

%------------------------------------------------
\begin{frame}{Techniques for Reducing Burnout}

\textbf{Core Idea: Thoughts influence feelings and behavior}



\begin{itemize}
    \item Cognitive Restructuring
    \begin{itemize}
        \item Identify and challenge negative thoughts
    \end{itemize}
    
    
    
    \item Thought Records
    \begin{itemize}
        \item Write situation $\rightarrow$ thought $\rightarrow$ alternative thought
    \end{itemize}
    
    
    
    \item Behavioral Activation
    \begin{itemize}
        \item Take small actions even without motivation
    \end{itemize}
    \end{itemize}

    \end{frame}
    \begin{frame}
    \begin{itemize}

    
    \item Task Breakdown
    \begin{itemize}
        \item Divide large tasks into manageable steps
    \end{itemize}
    
    
    
    \item Gradual Exposure
    \begin{itemize}
        \item Start with easier tasks, build up difficulty
    \end{itemize}
\end{itemize}

\end{frame}


%------------------------------------------------
\begin{frame}{Stress-Reframing Strategies}

\textbf{Core Idea: Stress depends on interpretation}



\begin{itemize}
    \item Stress-as-Enhancing Mindset
    \begin{itemize}
        \item Stress can improve focus and performance
    \end{itemize}
    
    
    
    \item Challenge vs. Threat
    \begin{itemize}
        \item View tasks as opportunities, not dangers
    \end{itemize}
    
    
    
    \item Normalization
    \begin{itemize}
        \item Burnout is common at the end of the semester
    \end{itemize}
    
    
    
    \item Process Focus
    \begin{itemize}
        \item Focus on effort (e.g., complete 5 problems)
    \end{itemize}
    
    
    
    \item Growth Framing
    \begin{itemize}
        \item Replace “I can’t” with “I’m still learning”
    \end{itemize}
\end{itemize}

\end{frame}
\begin{frame}{What Works Best?}
\begin{itemize}
    \item Strongest effects:
    \begin{itemize}
        \item Structured programs (not one-time activities)
        \item Repeated practice over time
    \end{itemize}
\end{itemize}



\begin{itemize}
    \item Less effective:
    \begin{itemize}
        \item One-off workshops
        \item Passive information delivery
    \end{itemize}
\end{itemize}

\end{frame}

\begin{frame}{Mechanisms}
Why do interventions work?



\begin{itemize}
    \item Improve coping strategies
    \item Increase sense of control
    \item Reduce cognitive overload
    \item Rebuild engagement
\end{itemize}



\begin{block}{Key Idea}
Burnout is reduced when students feel capable and supported
\end{block}

\end{frame}

%------------------------------------------------
\section{Pedagogical Implications}

\begin{frame}{Implications for Teaching}
\begin{itemize}
    \item Build structure into our remaining lessons 
    \item Use frequent, low-stakes activities
    \item Emphasize process over performance
\end{itemize}



\begin{itemize}
    \item Normalize struggle and fatigue
    \item Provide clear pathways for progress
\end{itemize}

\end{frame}

%------------------------------------------------

\begin{frame}{What Instructors Can Do}
\begin{itemize}
    \item Break large assignments into milestones
    \item Include short reflection activities
    \item Encourage peer interaction
\end{itemize}



\begin{itemize}
    \item Integrate small wellness practices
    \item Provide timely feedback
\end{itemize}

\end{frame}

%------------------------------------------------

\begin{frame}{End-of-Semester Strategies}
\begin{itemize}
    \item Focus on completion over perfection
    \item Highlight progress already made
    \item Reduce unnecessary cognitive load
\end{itemize}



\begin{block}{Key Pedagogical Shift}
Help students "get over the hump" rather than optimize performance
\end{block}

\end{frame}

%------------------------------------------------
\section{Discussion}

\begin{frame}{Discussion Questions}
\begin{itemize}
    \item Which type of intervention aligns most with your teaching?
    \item What barriers prevent implementing these strategies?
    \item How can we build ongoing support rather than one-time fixes?
\end{itemize}
\end{frame}

%------------------------------------------------

\begin{frame}{Discussion Questions (Continued)}
\begin{itemize}
    \item How do we balance rigor with burnout prevention?
    \item What small change could you implement next semester?
    \item How can departments (not just individuals) support this work?
\end{itemize}
\end{frame}

%------------------------------------------------

\begin{frame}{Takeaways}
\begin{itemize}
    \item Burnout is common and measurable
    \item Interventions work — especially when sustained
    \item Structure + support + engagement are key
\end{itemize}



\begin{block}{Final Thought}
Reducing burnout is not lowering standards — it is enabling students to meet them
\end{block}

\end{frame}

%------------------------------------------------

\end{document}


